% !TeX root = ../main.tex
% Add the above to each chapter to make compiling the PDF easier in some editors.

\chapter{Evaluation}\label{chapter:evaluation}
\section{Correctness}
We verify the correctness of our implementation by using the LLVM test suite and by running the SPEC CPU2017 integer benchmarks. Except for some LLVM test cases that contain unsupported instructions, all tests pass sucessfully on AArch64 and x86-64. Which we assume is a good proxy for the correctness of our implementation.

\section{Performance}

\subsection{Setup}
We evaluate the performance of our implementation by measuring the back-end
compile-time and the run-time performance of LLVM-IR generated by Clang. We use code produced
with -O1, where variables are in SSA form. We
compare our implementation against the original version of TPDE and against both the LLVM-O0 and LLVM-O1 back-end, the former two focusing on compile-time and the latter being a optimizing back-end. 
As benchmark programs, we use the SPEC CPU2017 integer benchmarks on x86-64 and AArch64.
Our x86-64 machine is an Amd Ryzen 5 5600X with 32 GiB of RAM running Linux \todo{version}; Our
AArch64 machine is an Apple M1 with 16 GiB of memory, using only the four performance cores,
running Asahi Linux \todo{version}. We use Clang/LLVM version 20.1.8.

